% Generated draft: requires mathematical and source review.
\documentclass[11pt]{article}
\usepackage[margin=25mm]{geometry}
\usepackage{fontspec}
\setmainfont{Latin Modern Roman}
\usepackage{amsmath,amssymb,mathtools,unicode-math}
\setmathfont{Latin Modern Math}
\usepackage{enumitem,needspace,xurl,hyperref}
\setlength{\parindent}{0pt}
\setlength{\parskip}{5pt}
\setlength{\emergencystretch}{3em}
\newcommand{\N}{\mathbb{N}}
\newcommand{\Z}{\mathbb{Z}}
\newcommand{\Q}{\mathbb{Q}}
\newcommand{\R}{\mathbb{R}}
\newcommand{\C}{\mathbb{C}}
\newcommand{\F}{\mathbb{F}}
\newcommand{\A}{\mathbb{A}}
\newcommand{\PP}{\mathbb{P}}
\newcommand{\E}{\mathbb{E}}
\newcommand{\eps}{\varepsilon}
\newcommand{\dd}{\,\mathrm{d}}
\begin{document}
\section*{UnsolvedMath ID: 30000769}
\textbf{Source problem code:} catalogue entry
\subsection{Definitions and mathematical statement}
For $n\geq3$, $M_{0,n}$ is the moduli space of smooth genus-zero curves with $n$ distinct ordered marked points, modulo isomorphism. Let $\mathcal P$ be the source's class of polygon pairs representing periods on these spaces. In the usual rational-algebra interpretation, let $A$ be the free commutative $\Q$-algebra on the corresponding formal symbols, and let $I$ be the ideal generated by the three specified relation families:
\begin{enumerate}
\item invariance under the symmetric-group relabeling;
\item one-point shuffle relations;
\item product-map relations.
\end{enumerate}
The proposed polygon algebra is $FC=A/I$. Let $\mathcal Z_{\mathrm{form}}$ denote the formal multiple-zeta-value algebra intended by the report. The conjectural target is an isomorphism
\[
FC\cong\mathcal Z_{\mathrm{form}},
\]
with the intended correspondence of polygon-period symbols and formal multiple-zeta values. The task is about the given period constructions, not an arbitrary abstract isomorphism of unrelated algebras.

The source also calls this equivalent to completeness of the three relation families among the corresponding periods. If these are numerical period relations, and $\operatorname{ev}:A\to\C$ sends each admissible symbol to its period, completeness means
\[
\ker(\operatorname{ev})=I.
\]
A formal-algebra isomorphism alone need not imply that numerical equality: the evaluation map from formal multiple-zeta values may itself have a kernel. Thus ``all relations'' must distinguish relations in the intended formal presentation from actual numerical period identities. The extract gives no presentation of $\mathcal Z_{\mathrm{form}}$, formulas for the three generators of $I$, comparison map, or polygon admissibility and regularization conventions. Their missing definitions and the asserted equivalence remain explicit rather than being filled by a chosen double-shuffle or associator presentation.
\subsection{Short English statement}
Ask whether the polygon-pair algebra modulo relabeling, one-point shuffle and product-map relations is the intended formal multiple-zeta algebra. The source also claims these relations are complete; formal identities and numerical period relations must be distinguished before that equivalence is precise.
\subsection{Sources}
\begin{itemize}
\item Supplied report: \emph{Arithmetic and Differential Galois Groups} (2008), Oberwolfach Reports, DOI 10.4171/owr/2007/26, pp. 1443--1520; proposed year 2007. \url{https://doi.org/10.4171/owr/2007/26}.
\item Local UnsolvedMath registry, permanent ID 30000769, catalogue entry; created 25 August 2026, updated 16 September 2026. Its corrected\_verified note says the report was checked directly and all three relation families and the claimed equivalence were restored. The supplied 21 August 2026 literature triage labels the comparison conjectural and open, noting verification only in small weights. No additional external reference is supplied; the dated assessment has not been independently verified here.
\end{itemize}
\end{document}
